\chapter{Autenticazione e gestione delle identità}
\label{chap:autenticazione}

\section{Ciclo di vita della richiesta}
\label{sec:ciclo-richiesta}

Una richiesta di registrazione o di login attraversa tre servizi prima di diventare un record:
Entry-Hub, Security-Switch e Database-Vault. Tutti e tre applicano le stesse verifiche sull'input,
ognuno indipendentemente dagli altri.

La validazione avviene in due fasi. La prima riguarda il corpo della richiesta: il payload viene
letto fino a un limite di 2048 byte e rifiutato se lo supera, senza accumulare in memoria dati
controllati da chi ha inviato la richiesta, e il decoder JSON respinge qualunque campo inatteso.

La seconda fase riguarda i singoli campi \cite{ramusb-code-eh-f-04}. L'email deve essere un
indirizzo singolo e sintatticamente valido secondo RFC 5322 \cite{rfc5322}; la password deve essere
lunga tra 8 e 128 caratteri e usare almeno tre categorie tra minuscole, maiuscole, cifre e simboli.
In registrazione, inoltre, la chiave pubblica SSH deve essere ben formata.

Se una verifica fallisce, la risposta è HTTP 400 senza indicare quale controllo non sia passato, il
log registra il problema senza identificare l'utente, e la richiesta non prosegue verso l'interno. 
Il messaggio resta generico perché un errore dettagliato aiuterebbe un attaccante più di quanto aiuti un
utente in buona fede, che le stesse verifiche le ha già ricevute dal proprio client.

I tre servizi condividono la stessa libreria di validazione, quindi ripetere il controllo non
introduce diversità algoritmica, ma garantisce che nessun servizio dia per scontato che il
precedente abbia fatto il proprio lavoro. In questo modo, se a Database-Vault dovesse arrivare una
richiesta proveniente da un servizio compromesso, la validazione verrebbe fatta comunque.

Gli errori interni vengono infine tradotti in codici HTTP. La \autoref{tab:codici-http} riassume i codici 
che raggiungono l'utente, attribuendo ciascun codice a chi lo emette.

Entry-Hub costruisce soltanto gli errori che nascono da sé stesso: una validazione fallita e una
chiamata a Security-Switch andata male. Tutto il resto lo rilancia al client così come gli arriva
dall'interno, senza modificarlo. Per questo motivo, un 401 o un 409 raggiungono l'utente pur non
essendo mai prodotti da Entry-Hub.

\begin{table}[H]
\centering
\caption{Codici HTTP restituiti al client e loro significato nel sistema.}
\label{tab:codici-http}
\begin{tabular}{@{}l p{\dimexpr0.58\linewidth\relax} l@{}}
\hline
\textbf{Codice} & \textbf{Significato} & \textbf{Emesso da} \\
\hline
400 & Validazione fallita, senza indicare quale controllo non sia passato & entrambi \\
401 & Credenziali non valide, non specifica se email inesistente o password errata & Security-Switch \\
403 & Rifiuto esplicito di Network-Manager & Security-Switch \\
409 & Email o chiave SSH già registrate, non specifica quale delle due & Security-Switch \\
500 & Errore interno non altrimenti classificato & entrambi \\
502 & Servizio a valle irraggiungibile o con risposta inattesa & entrambi \\
503 & Chiamata a Security-Switch scaduta & Entry-Hub \\
504 & Chiamata a Database-Vault o Network-Manager scaduta & Security-Switch \\
\hline
\end{tabular}
\end{table}

Tra i codici che i due servizi costruiscono per conto proprio ci sono due differenze: 
la prima è il 403, che solo Security-Switch può emettere, perché è l'unico dei due a poter ricevere un 
rifiuto esplicito da un servizio interno. La seconda riguarda le chiamate scadute: Entry-Hub risponde 503, 
Security-Switch 504.

In tutti i casi il corpo della risposta è sanificato, e il dettaglio dell'errore resta nei log
interni.

\section{Cifratura e hashing delle credenziali}

Database-Vault tratta i dati che riceve in tre modi diversi:

\begin{itemize}
    \item la chiave SSH è pubblica, quindi non necessita di segretezza;
    \item l'email è insieme una chiave di ricerca e un dato personale da proteggere;
    \item la password non deve essere recuperabile, nemmeno da chi amministra il database.
\end{itemize}

L'email viene normalizzata in minuscolo e ridotta al proprio digest SHA-256, che diventa la chiave primaria 
del record. La normalizzazione garantisce che lo stesso indirizzo scritto con lettere diverse al login 
produca la stessa chiave della registrazione.
Questo hash è deterministico per necessità, perché deve permettere il lookup. Il limite fondamentale di 
questo approccio, insieme alla sua mitigazione, è discusso nel \autoref{chap:sicurezza}.

Accanto al digest che fa da chiave di ricerca, il record conserva una copia cifrata, così che 
l'indirizzo resti recuperabile. Da una master key e da un salt casuale di 16 byte, generato per quel record, 
viene derivata con HKDF-SHA256 \cite{rfc5869} una chiave dedicata, con cui l'email è cifrata in
AES-256-GCM \cite{nist-sp-800-38d} usando un nonce casuale di 12 byte. Salt, nonce e testo cifrato
vengono salvati nel record, mentre la chiave derivata viene azzerata appena il cifrario l'ha consumata.
Derivare una chiave per ogni record, invece di cifrare tutto con la master key, riduce la quantità
di dati protetti da una singola chiave.

\Needspace{22\baselineskip}
\begin{lstlisting}[caption={Derivazione della chiave per record in \texttt{newGCM}, da
\texttt{encryption.go} \cite{ramusb-code-hkdf-newgcm}.}, label={lst:hkdf}]
func newGCM(masterKey, salt []byte) (cipher.AEAD, error) {
	derivedKey := make([]byte, derivedKeySize)
	kdf := hkdf.New(sha256.New, masterKey, salt, []byte(hkdfInfo))
	if _, err := io.ReadFull(kdf, derivedKey); err != nil {
		return nil, fmt.Errorf("encryption: derive key: %w", err)
	}

	block, err := aes.NewCipher(derivedKey)
	for i := range derivedKey {
		derivedKey[i] = 0
	}
	if err != nil {
		return nil, fmt.Errorf("encryption: new AES cipher: %w", err)
	}

	gcm, err := cipher.NewGCM(block)
	if err != nil {
		return nil, fmt.Errorf("encryption: new GCM: %w", err)
	}

	return gcm, nil
}
\end{lstlisting}


La password viene invece concatenata a un pepper, letto da una variabile d'ambiente, e passata ad Argon2id 
\cite{rfc9106} insieme a un salt casuale di 16 byte generato sul momento per quel record.

OWASP \cite{owasp-password-storage} propone diverse configurazioni equivalenti tra loro, che
scambiano memoria e numero di passaggi. RAM-USB parte da quella con più memoria e ne raddoppia i
passaggi, pagando quindi a ogni verifica il doppio del lavoro che OWASP considera sufficiente. È un
costo che rende impraticabile recuperare le password da un database rubato. Il pepper complica
ulteriormente un attacco offline: a differenza del salt, che è salvato nel record e impedisce
soltanto che un tentativo valga per più utenti insieme, il pepper nel database non compare: in sua
assenza nessun tentativo è calcolabile, e una copia del database non è quindi sufficiente a condurre
l'attacco.

Il risultato viene salvato come stringa autodescrittiva in formato PHC, la convenzione di codifica
nata dalla Password Hashing Competition. Ogni record porta quindi con sé i parametri con cui è stato creato, 
e la verifica al login li usa al posto delle costanti correnti del servizio. 
In questo modo, cambiare i parametri per i nuovi utenti non invalida i record già esistenti. 

La \autoref{tab:phc} elenca i campi di quella stringa, nell'ordine in cui compaiono nel record.

\begin{table}[H]
\centering
\caption{Campi della stringa PHC in cui Database-Vault salva l'hash della password.}
\label{tab:phc}
\begin{tabular}{@{}l p{\dimexpr0.70\linewidth\relax}@{}}
\hline
\textbf{Campo} & \textbf{Significato} \\
\hline
\texttt{argon2id} & Variante dell'algoritmo \\
\texttt{v=19} & Versione di Argon2: il byte \texttt{0x13} fissato da RFC 9106, stampato in decimale \\
\texttt{m=47104} & Memoria di lavoro in KiB, cioè 46 MiB \\
\texttt{t=2} & Numero di passaggi sulla memoria \\
\texttt{p=1} & Grado di parallelismo \\
salt & Salt casuale di 16 byte, generato per quel record \\
digest & Hash di 32 byte prodotto da Argon2id \\
\hline
\end{tabular}
\end{table}

Master key e pepper risiedono entrambi in variabili d'ambiente, senza una procedura di backup né di
rotazione. Il rischio è discusso nel \autoref{chap:conclusioni}.


Lo \autoref{sch:credenziali} riassume le trasformazioni appena descritte, con i simboli seguenti:

\begin{itemize}
    \item $e$: l'email ricevuta;
    \item $p$: la password ricevuta;
    \item $k_{SSH}$: la chiave pubblica SSH;
    \item $\pi$: il pepper;
    \item $MK$: la master key;
    \item $s$, $n$, $s_p$: rispettivamente il salt di HKDF, il nonce di AES-GCM e il salt di
    Argon2id, generati per quel record;
    \item $\stackrel{R}{\leftarrow}$: un'estrazione casuale uniforme dall'insieme indicato.
\end{itemize}

\begin{schema}
\centering
$\begin{array}{rl}
1. & s,\, s_p \stackrel{R}{\leftarrow} \{0,1\}^{128}, \quad n \stackrel{R}{\leftarrow} \{0,1\}^{96} \\[4pt]
2. & k_h = \mathrm{SHA\mbox{-}256}(\mathrm{lower}(e)) \\[4pt]
3. & k = \mathrm{HKDF\mbox{-}SHA256}(MK,\, s) \\[4pt]
4. & c = \mathrm{AES\mbox{-}GCM}_{k}(n,\, e) \\[4pt]
5. & h = \mathrm{Argon2id}(p \, \| \, \pi,\, s_p) \\[4pt]
6. & r = (k_h,\, s,\, n,\, c,\, \mathrm{PHC}(h,\, s_p),\, k_{SSH})
\end{array}$
\caption{Trasformazioni applicate alle credenziali prima della scrittura del record.}
\label{sch:credenziali}
\end{schema}

\subsection{Robustezza della password}
\label{sec:entropia}

Le regole di validazione della \autoref{sec:ciclo-richiesta} delimitano lo spazio delle password
ammesse, e da quello spazio si ricava, nel caso migliore, quanta incertezza la policy possa
garantire a chi tenta di indovinare una password, ossia quante possibilità debba considerare, non
sapendo quale sia stata scelta.

Quanto valga quell'incertezza dipende però da come la password è stata scelta. Nel caso migliore, in
cui è prodotta da un generatore pseudo-casuale che pesca dall'intero alfabeto ammesso, ogni
password ammessa è equiprobabile e l'incertezza
è quindi massima. Nel caso peggiore, in cui è costruita da una persona, alcune password sono
ben più probabili di altre, ad esempio, \texttt{Password1!} viene scelta molto più spesso di
una sequenza casuale della stessa lunghezza. Il valore calcolato qui è quello del caso migliore, cioè un
limite superiore.

In questo caso, quindi, il calcolo si riduce a contare i valori distinti che la policy ammette, a partire 
dalla cardinalità delle quattro categorie:

\[
\begin{array}{rcll}
n_1 & = & 26  & \mbox{minuscole} \\
n_2 & = & 26  & \mbox{maiuscole} \\
n_3 & = & 10  & \mbox{cifre} \\
n_4 & = & 33  & \mbox{simboli} \\[2pt]
N   & = & \sum_i n_i = 95 & \mbox{caratteri disponibili in totale}
\end{array}
\]

Il vincolo di almeno tre categorie esclude le stringhe che ne usano al più due, quindi per contare le
password valide conviene contare quelle che non lo sono: si parte da tutte le $N^L$ stringhe
possibili e si sottraggono quelle che usano una sola categoria e quelle che ne usano esattamente due.

Il primo gruppo è immediato, perché le stringhe scritte con la sola categoria $i$ sono $n_i^L$. Per
il secondo non basta sommare $(n_i + n_j)^L$ su tutte le coppie, perché quel valore conta le stringhe
scritte con quei due soli alfabeti e comprende quindi anche quelle di sola $i$ e di sola $j$:
sottraendo $n_i^L$ e $n_j^L$ restano le stringhe che usano davvero entrambe le categorie. Per il
principio di inclusione-esclusione, il numero di password valide di lunghezza $L$ vale quindi

\[
V(L) = N^L - \sum_{i=1}^{4} n_i^L
       - \sum_{i<j} \left[ (n_i + n_j)^L - n_i^L - n_j^L \right].
\]

Un attaccante privo di qualsiasi informazione deve considerare tutti i $V(L)$ candidati ammessi.
Quel numero si esprime abitualmente in bit, cioè come esponente di 2: $b$ bit corrispondono a $2^b$
candidati. Se la password è estratta uniformemente, l'incertezza dell'attaccante vale perciò 
$\log_2 V(L)$ bit. 

Per la lunghezza minima ammessa, $L = 8$, si ottiene $V(8) = 6{,}27 \cdot 10^{15}$, cioè
$\log_2 V(8) \approx 52{,}48$ bit.

Il confronto con lo spazio non vincolato mostra quanto pesi ciascuna delle due regole. Senza il
vincolo sulle categorie sarebbero ammesse tutte le stringhe di otto caratteri sull'alfabeto, cioè
$N^8 = 95^8 = 6{,}63 \cdot 10^{15}$, cioè 52,56 bit. La suddetta regola quindi riduce i bit di 0,08. 
Aggiungere un carattere in più, invece, aumenta i bit di entropia di $\log_2 95 \approx 6{,}57$ bit per ogni
carattere aggiunto.

Da queste cifre potrebbe sembrare che il vincolo semplifichi il lavoro di un malintenzionato invece
di contrastarlo, ed è vero, per quanto in misura trascurabile, nel caso in cui le password siano create da 
generatori pseudo-casuali. 
Lo scopo di quella regola però non è aumentare l'entropia, bensì vincolare la scelta di
una persona, che tende a fermarsi sulle sole minuscole o su lettere e cifre. Quanto valga in quel
ruolo non è ricavabile da questo calcolo.

Se le due regole debbano coesistere, o se convenga sostituire il vincolo sulle categorie con una
lunghezza minima maggiore, resta quindi una questione aperta, ripresa tra gli sviluppi futuri nel
\autoref{chap:conclusioni}.


Quanto quello spazio resista dipende però anche dal costo di ogni tentativo, cioè dai parametri di
Argon2id: il \autoref{chap:sicurezza} quantifica quel costo dal punto di vista di chi attacca.

\section{Registrazione}

Quando la richiesta raggiunge Database-Vault ed è stata rivalidata, il servizio calcola le tre
trasformazioni della sezione precedente e prova a scrivere il record in una singola transazione
atomica \cite{ramusb-code-saveuser}. Due vincoli di unicità proteggono la tabella, uno
sull'hash dell'email e uno sulla chiave pubblica SSH: se uno dei due viene violato, la registrazione
è rifiutata con HTTP 409, senza dire quale dei due abbia causato il conflitto. Quel codice
resta però un'informazione: dice a un mittente non autenticato che uno dei due valori inviati era
già presente, e chi invia una chiave SSH appena generata sa che l'unico altro candidato è l'email.
Il \autoref{chap:sicurezza} analizza questo problema e presenta una possibile mitigazione.

A questo punto il record esiste, ma l'utente non ha ancora uno spazio in cui scrivere.
Database-Vault genera allora un nome utente POSIX nella forma
\texttt{user\textless{}xxxxxx\textgreater{}}, dove le sei cifre sono caratteri casuali di un
alfabeto base-36, e chiede a Storage-Service di creare l'account corrispondente. Come
Storage-Service lo crei e lo isoli è oggetto del \autoref{chap:storage}.

Se quella creazione fallisce, il record appena scritto viene eliminato. Se anche il rollback fallisce, 
il servizio lo segnala con un errore dedicato.

Quando invece entrambi i passi riescono, Database-Vault comunica a Security-Switch che la
registrazione è andata a buon fine, e quest'ultimo procede con l'orchestrazione descritta
nella sezione conclusiva di questo capitolo.

\Needspace{12\baselineskip}
\begin{lstlisting}[caption={Rollback compensativo dopo un fallimento di Storage-Service, da
\texttt{registration.go} \cite{ramusb-code-registration-rollback}.}, label={lst:rollback}]
		rollbackCtx, cancel := context.WithTimeout(context.WithoutCancel(ctx), rollbackTimeout)
		defer cancel()

		if delErr := store.DeleteUser(rollbackCtx, input.EmailHash); delErr != nil {
			return Result{
				Outcome: OutcomeFailed,
				Err:     fmt.Errorf("%w: posix creation error: %w; delete error: %w", ErrOrphanedRecord, err, delErr),
			}
		}
\end{lstlisting}

La \autoref{fig:uc01} riassume l'intera sequenza, dalla richiesta del client fino alla risposta. Gli
ultimi due scambi, quelli verso Network-Manager, sono oggetto della sezione conclusiva di questo
capitolo.

\begin{figure}[H]
    \centering
    \includegraphics[width=\textwidth]{figures/03-usecases-sequence-uc01-registration.pdf}
    \caption{Sequenza di registrazione (UC-01).}
    \label{fig:uc01}
\end{figure}

\section{Login}

Durante il login, Database-Vault ricalcola il digest SHA-256 dell'email ricevuta e con quello cerca il 
record.

Salt e parametri di costo vengono estratti dalla stringa PHC del record. Argon2id viene ricalcolato con 
quei valori e con il pepper, e il digest ottenuto è confrontato con quello salvato tramite un confronto a 
tempo costante, che non lascia quindi dedurre quanti byte iniziali dei due valori coincidano.

Un'email inesistente e una password sbagliata devono produrre la stessa risposta, e non devono
essere distinguibili nemmeno nei log. Le due strade convergono quindi su un unico errore
sentinella, e il contenuto dell'errore originale viene scartato invece che riportato
\cite{ramusb-code-login}.

Rendere identici la risposta e il log non basta però a rendere indistinguibili i due casi, perché
resta una terza differenza osservabile: il tempo. Argon2id, con i parametri della sezione
precedente, costa decine di millisecondi, e se venisse eseguito solo quando un record esiste
davvero, un attaccante potrebbe distinguere le due situazioni misurando la latenza della risposta. 
La soluzione adottata è spendere lo stesso tempo anche quando non c'è nulla da verificare.

\Needspace{6\baselineskip}
\begin{lstlisting}[caption={Verifica fittizia usata sui percorsi di rifiuto, da \texttt{hash.go}
\cite{ramusb-code-verify-dummy}.}, label={lst:verifydummy}]
func VerifyDummy(password, pepper []byte) {
	_, _ = VerifyPassword(password, pepper, dummyHash())
}
\end{lstlisting}

La funzione ricalcola Argon2id contro un hash fittizio, costruito una sola volta per processo con
gli stessi parametri di costo dei record veri, e ne scarta il risultato. Viene chiamata su ogni
percorso di rifiuto, compreso quello in cui è il lookup stesso a fallire, così nemmeno un guasto
dell'infrastruttura diventa facilmente misurabile dall'esterno.

Ogni incertezza porta allo stesso esito. Se la verifica non può concludersi, per esempio perché
l'hash salvato è malformato, la risposta resta 401, secondo il principio
fail-secure \cite{owasp-fail-securely}. Nei log quel caso resta distinguibile, perché a un operatore
serve sapere che si tratta di un record corrotto, di un utente distratto o di un vero e proprio attacco.

La \autoref{fig:uc02} riassume la sequenza completa di login, compreso il ramo in cui le due cause
di rifiuto convergono sulla stessa risposta.

\begin{figure}[H]
    \centering
    \includegraphics[width=\textwidth]{figures/04-usecases-sequence-uc02-login.pdf}
    \caption{Sequenza di autenticazione (UC-02).}
    \label{fig:uc02}
\end{figure}

\section{Orchestrazione post-autenticazione}

Autenticare un utente non gli dà comunque ancora accesso a niente. L'accesso, in questo sistema, è
una proprietà della rete. Security-Switch richiede a Network-Manager la modifica delle regole di rete
in due momenti diversi.

Dopo una registrazione riuscita chiede la creazione dell'identità dell'utente sulla rete mesh e di
una pre-auth key, che risale la catena fino al client dentro la risposta 201 e che permette al suo
dispositivo di entrare nella rete privata.

Dopo un login riuscito chiede invece un accesso temporaneo di 12 ore verso Storage-Service.

Entrambi i passaggi concedono però raggiungibilità, non lettura: la registrazione mette il
dispositivo nella rete, il login gli permette di raggiungere Storage-Service, ma ad aprire i backup è
la chiave privata SSH, che non lascia mai il dispositivo dell'utente. 
Conoscere le credenziali dell'utente non è pertanto condizione sufficiente ad accedere ai suoi backup.

Come Network-Manager crei l'identità mesh, applichi il tag ACL e ne gestisca la scadenza è oggetto
del \autoref{chap:rete}.